\section{Conclusion and Future Work}

In this report, we presented LongCat-Video-Avatar 1.5, an open-source framework for audio-driven video generation tailored for practical applications. By focusing on empirical optimization and production readiness, we aim to narrow the gap between academic research prototypes and industrial deployments. The integration of the Whisper-large audio encoder, combined with our comprehensive data curation pipeline and scaled multi-stage training recipe, enables the model to achieve highly precise lip-synchronization, full-body temporal stability, and reliable identity consistency in long-horizon video generation.

Furthermore, our model demonstrates robust adaptability to complex real-world conditions, such as multi-person conversations, object handling, and stylized domains (e.g., anime and animals). Through the application of Group-Relative Policy Optimization (GRPO) for human preference alignment and Distribution Matching Distillation (DMD) for inference acceleration, we developed an efficient 8-NFE inference pipeline that effectively balances generation speed and visual fidelity. Extensive human evaluations across diverse scenarios indicate that LongCat-Video-Avatar 1.5 achieves highly competitive performance in naturalness, stability, and overall visual realism when compared to existing closed-source systems like OmniHuman 1.5 and HeyGen.

\paragraph{Future work.} Regarding future work, current virtual human generation models still have significant room for improvement in physical plausibility and fine-grained audio-visual synchronization. Furthermore, to maintain long-term identity consistency, existing methods often over-rely on fixed reference frames. This reliance inevitably leads to motion repetition and unnatural camera transitions constrained by the reference views. Therefore, developing a truly unbounded, infinite-length video generation framework that inherently preserves identity without rigid dependence on static reference frames remains a critical direction for future research. 